\section{The endpoint problem and the complete source}
\label{sec:intro}

The purpose of this paper is to assemble the analytic results of the
relative-charge investigation into a single argument with explicit
hypotheses. The question motivating the investigation is whether the
unchanged theta source supplies enough signed information to exclude a
multiple zero at time zero, conditional on the Riemann hypothesis (RH).
The results below supply geometric windows, exact response identities,
and quantitative estimates to which such information could be applied.
They also identify limitations of several simpler substitutes.

\subsection{Normalization and analytic bounds}

Throughout the theta application, put
\begin{align}
 H(t,z)&=\int_0^\infty \Phi(u)e^{tu^2}\cos(zu)\,du,\label{i:source}\\
 \Phi(u)&=\sum_{n\geq1}\phi_n(u),\qquad \lambda_n=\pi n^2,\notag\\
 \phi_n(u)&=(2\lambda_n^2e^{9u}-3\lambda_ne^{5u})
                  e^{-\lambda_ne^{4u}}.\label{i:phi}
\end{align}
This is the standard normalization used in \cite{polymath}:
\begin{equation}
 H(0,z)=\frac18\xi\!\left(\frac12+\frac{iz}{2}\right),\qquad
 \xi(s)=\frac{s(s-1)}2\pi^{-s/2}\Gamma(s/2)\zeta(s).
 \label{i:xi}
\end{equation}
The Fourier representation \eqref{i:xi} and the classical real-zero
threshold theorem are the analytic number-theory inputs; none of the
local proofs below uses a zero table or RH unless it is explicitly stated.

For $u\geq0$ every summand in \eqref{i:phi} is positive, since
$2\lambda_ne^{4u}-3>0$. The series and its derivatives converge
uniformly on compact real $u$ intervals. For large $u$ the summands are
bounded by a polynomial in $n,e^u$ times $e^{-\pi n^2e^{4u}}$;
factoring out, for example, $e^{-(\pi/2)e^{4u}}$ leaves a uniformly
summable series. Consequently $\Phi$ and each of its fixed-order
derivatives decay faster than $e^{-cu^2-du}$ for any fixed real $c,d$.
Differentiation under the integral is therefore justified locally
uniformly in $(t,z)\in\mathbb C^2$. In particular,
\begin{equation}
 H_t=-H_{zz},\qquad
 H(t,-z)=H(t,z),\qquad
 H(\overline t,\overline z)=\overline{H(t,z)}.
 \label{i:heat}
\end{equation}
Here and below subscripts denote derivatives; the time parameter is
written as an argument to avoid confusing it with a derivative.

For $T\in\mathbb R$, $R\geq0$ and nonnegative integer $j$, define the
\emph{complete moment}
\begin{equation}
 M_j(T,R)=\int_0^\infty \Phi(u)u^j e^{Tu^2}\cosh(Ru)\,du.
 \label{i:moments}
\end{equation}
Since $|\cos(x+iy)|$ and $|\sin(x+iy)|$ are at most $\cosh y$,
\begin{equation}
 |\partial_t^a\partial_z^bH(t,z)|\leq M_{2a+b}(T,R)
 \quad(\operatorname{Re}t\leq T,\ |\operatorname{Im}z|\leq R).
 \label{i:jet-bound}
\end{equation}
These bounds are independent of $\operatorname{Re}z$. ``Complete'' means
that every source mode and the entire integration tail are included.
For a signed residual source $\psi$, the same assertion holds with
$|\psi|$ in place of $\Phi$.

\subsection{What RH supplies, and what it leaves open}

The de Bruijn--Newman theorem gives a finite real $\Lambda$ for which
all zeros of $H(t,\cdot)$ are real exactly when $t\geq\Lambda$.
In this normalization RH is equivalent to $\Lambda\leq0$;
Rodgers and Tao proved $\Lambda\geq0$ \cite{rodgers-tao}.
Thus under RH the zeros are real for all $t\geq0$. They are simple for
$t>0$: if a real zero had multiplicity $m\geq2$ at $t_0>0$, the local
Hermite splitting proved in Proposition~\ref{f:hermite}, applied at
$t=t_0-\varepsilon$, would give nonreal zeros for sufficiently small
$\varepsilon>0$. At least two Hermite roots are nonzero, and their
imaginary displacement has order $\sqrt{\varepsilon}$, whereas the
error is $O(\varepsilon)$. This contradicts reality at that time.
The argument does not rule out a multiple zero at $t=0$.

A common negative time at which \emph{all} theta zeros are real is
unavailable, even under RH, since it would contradict $\Lambda\geq0$.
This restriction does not prohibit different negative-time neighborhoods
for different endpoint heights. The quantitative windows below are
explicitly allowed to depend on the hypothetical event.

The local charge at an endpoint of multiplicity $m$ will be
$q=\binom m2$. It relates directly to the multiplicity excess studied
in the companion heat-motion investigation \cite{bw10,assessment07}.
Under RH let $\rho=\tfrac12+i\gamma$ range over distinct nontrivial
zeros of both ordinate signs, with multiplicity $m_\rho$. For every
finite symmetric cutoff $Y$, define
\begin{equation}
 K_Y=\sum_{|\gamma|\leq Y}\frac{m_\rho^2}{|\rho|^2},\qquad
 Z_Y=\sum_{|\gamma|\leq Y}\frac{m_\rho}{|\rho|^2}.
 \label{i:weighted}
\end{equation}
Termwise algebra gives
\begin{equation}
 K_Y-Z_Y=2\sum_{|\gamma|\leq Y}\frac{q_\rho}{|\rho|^2}.
 \label{i:bridge}
\end{equation}
When the two weighted sums converge, as in the companion framework,
the same identity holds without the cutoff. No convergence theorem is
needed for the local arguments in this paper. At the physical coordinate
$a=2\gamma$, the charge weight on the right is $8/(1+a^2)$.
Thus the local integer charge does not produce a positive uniform gap
for the weighted excess. An endpoint-local charge must also be
distinguished from the charge of a larger window, which might contain a
negative-time event.

\subsection{Degree, relative Chern number, and Hopf interpretation}

Let $\Omega\subset\mathbb C^2$ be a bounded domain with piecewise smooth
boundary, and suppose $G=(H,H_z)$ is nonzero on that boundary. Equip both
copies of $\mathbb C^2$ with the complex orientation, in the coordinate
order $(\operatorname{Re}t,\operatorname{Im}t,
\operatorname{Re}z,\operatorname{Im}z)$. The boundary charge is
\begin{equation}
 Q_\Omega=\deg\left(\frac{G}{\|G\|}:\partial\Omega\longrightarrow S^3\right).
 \label{i:degree}
\end{equation}
It is the relative top Chern number of the trivial complex rank-two
bundle with the specified nonvanishing boundary section: equivalently,
the total index of its interior zeros. For isolated holomorphic zeros
these indices are positive analytic intersection multiplicities.
Theorem~\ref{f:charge} supplies the discriminant computation of them.

For the product of a time rectangle and a spatial disk, rounding corners
makes the boundary an oriented $S^3$. Its usual Hopf interpretation
requires no additional charge formula. Let $\pi:S^3\to S^2$ be the Hopf
map, let $\omega$ have integral one on $S^2$, and choose a connection
form $\alpha$ with $d\alpha=\pi^*\omega$ and
$\int_{S^3}\alpha\wedge d\alpha=1$. For $g=G/\|G\|$ the Hopf invariant
of $\pi\circ g$ is
\[
 \int_{\partial\Omega}g^*\alpha\wedge d(g^*\alpha)
 =\int_{\partial\Omega}g^*(\alpha\wedge d\alpha)=\deg g.
\]
Here the integral formula defines the normalized Hopf invariant; its
independence of the chosen primitive follows by Stokes' theorem and
$H^1(S^3)=0$. It agrees with \eqref{i:degree}. The analytic estimates
in this paper use degree and discriminants directly.

Neither reality nor spatial parity supplies cancellation. Conjugating
both complex domain coordinates preserves real orientation, as does
conjugating both target coordinates. The change $z\mapsto-z$ and its
target action $(H,H_z)\mapsto(H,-H_z)$ are also orientation preserving.
The corresponding reflected events therefore have the same positive
local charge.

\subsection{Reading the results}

The first part proves proper cluster preparation, local charge, the
two-root cap theorem and explicit event-compatible windows. The second
part retains signs in source perturbations, proves quantitative event
continuation, and specializes the response to the actual canonical theta
tail. The many-root part keeps every root of a higher birth and shows
exactly where the two-root argument ceases to apply. Numerical evidence
is collected separately in Section~\ref{sec:support}. Appendix~\ref{app:reference}
records the holomorphic reference comparison and the boundary
interpolation contract used by earlier certificates.

All propositions are statements with the hypotheses displayed here.
In particular, a narrower source-response interval need not lie on one
side of zero, and a phase bound need not imply zero charge. This
manuscript consolidates internally reviewed research results; it has not
received independent specialist validation. It asserts no priority for
the standard complex-analysis and heat-root-motion ingredients.
